%────────────────────────────────────────────────────────────────────────────
% Template LaTeX per relazione in stile articolo scientifico
% Autore: [Tuo Nome]
% Data: \today
%────────────────────────────────────────────────────────────────────────────
\documentclass[11pt,a4paper]{article}

% ────────────────────────────────────────────────────────────
% Impaginazione e font
% ────────────────────────────────────────────────────────────
\usepackage[margin=2.5cm]{geometry}   % margini generosi
\usepackage{lmodern}                  % font Latin Modern (professionale)
\usepackage[T1]{fontenc}              % giusta codifica font

% ────────────────────────────────────────────────────────────
% Tipografia e microtipografia
% ────────────────────────────────────────────────────────────
\usepackage{microtype}                % migliore giustezza dei caratteri
\usepackage{setspace}                 % controllo spaziatura
\onehalfspacing                       % 1.5 interlinea (più leggibile)

% ────────────────────────────────────────────────────────────
% Matematica e simboli
% ────────────────────────────────────────────────────────────
\usepackage{amsmath,amssymb,amsfonts} % pacchetti AMS
\usepackage{bm}                       % grassetto per vettori/matrici
\usepackage{slashed}                  % \slashed{p} (utile in QCD)

% ────────────────────────────────────────────────────────────
% Figure, tabelle e codice
% ────────────────────────────────────────────────────────────
\usepackage{graphicx}                 % includere immagini
\usepackage{subcaption}               % sottocaption per subfigure
\usepackage{booktabs}                 % tabelle professionali
\usepackage{caption}                  % personalizzare caption
\usepackage{listings}                 % codice sorgente
\lstset{
  basicstyle=\ttfamily\small,
  breaklines=true,
  frame=single,
  captionpos=b
}

% ────────────────────────────────────────────────────────────
% Link e riferimenti incrociati
% ────────────────────────────────────────────────────────────
\usepackage[hidelinks]{hyperref}      % link senza colori

% ────────────────────────────────────────────────────────────
% Header & Footer
% ────────────────────────────────────────────────────────────
\usepackage{fancyhdr}
\pagestyle{fancy}
\fancyhf{}
\lhead{Relazione su QCD}
\rhead{[Tuo Nome]}
\cfoot{\thepage}

% ────────────────────────────────────────────────────────────
% Titolo, autori, abstract
% ────────────────────────────────────────────────────────────
\title{\textbf{Studio di QCD tramite Mathematica}}
\author{[Tuo Nome] \\ \small Università degli Studi di …}
\date{\today}

\begin{document}

\maketitle

\begin{abstract}
  Vogliamo usare Mathematica per generare ampiezze di Feynman, quindi modellizzare una situazione fisica (in questo caso di fisica delle particelle), introducendo un formalismo e creando notazioni.

  Per fare un modello dobbiamo scegliere prima di tutto il tipo di interazione, noi studieremo i processi $gg \rightarrow gg$ e $gg \rightarrow ggg$, quindi processi puramente di QCD.

  Prime di tutto implementiamo delle procedure per generare tutti i possibili grafi connessi collegando tutti i punti tra loro di loro. I punti esterni vanno tenuti fermi e poi disegnamo tutte le possibili linee interne.

  Una volta che ho tutti i grafi ho una base geometrica per generare qualcosa di fisico: regole di Feynman e conservazione del quadriimpulso di ciascun diagramma (la conservazione di tutti gli impulsi entranti deve fare zero, l'idea è di trovare un modo automatico per implementare la conservazione dell'impulso a ciascun vertice).

  Implementiamo una teoria di campo per il nostro modello, indicizzando i punti esterni con interi positivi e i punti interni con interi negativi. Le linee interne sono propagatori a due punti.

  Livello successivo: Lagrangiana, regole di Feynman.

\end{abstract}

% ────────────────────────────────────────────────────────────
% Struttura tipica: sezioni
% ────────────────────────────────────────────────────────────
\section{Introduzione}
% - Contestualizza QCD
% - Motivazioni del lavoro
% - Breve panoramica dei contenuti

\section{Formalismo Teorico}
% - Lagrangiana della QCD
% - Simmetrie e proprietà fondamentali
% - Quantizzazione e gauge fixing

\section{Implementazione in Mathematica}
% - Uso di pacchetti (FeynCalc, etc.)
% - Codice esemplificativo
\begin{lstlisting}[language=Mathematica,caption={Esempio di calcolo con FeynCalc}]
(* Setup *)
<< FeynCalc`
(* Definizione Lagrangiana *)
lag = GAD[\mu,\nu]^2/(4*g^2);
(* Calcolo di un propagatore *)
Prop = FAD[{q,0}];
\end{lstlisting}

\section{Risultati}

\section{Discussione}
% - Interpretazione dei risultati
% - Confronto con letteratura

\section{Conclusioni}
% - Punti chiave emersi
% - Prospettive future

\appendix
\section{Dettagli di calcolo}
% - Espandi formule troppo lunghe
% - Codice completo se serve

\bibliographystyle{unsrt}
\bibliography{bibliografia}  % crea bibliografia in file .bib

\end{document}
